\newpage
\section{第21章：递归类型的元理论}

\begin{taplauthorsolutionentrybox}
\phantomsection
\label{sol:author:21.1.7}
\textbf{21.1.7 解答}

\[
\begin{array}{c@{\quad=\quad}c@{\qquad}c@{\quad=\quad}c}
E_2(\varnothing)&\{a\}&E_2(\{a,b\})&\{a,c\}\\
E_2(\{a\})&\{a\}&E_2(\{a,c\})&\{a,b\}\\
E_2(\{b\})&\{a\}&E_2(\{b,c\})&\{a,b\}\\
E_2(\{c\})&\{a,b\}&E_2(\{a,b,c\})&\{a,b,c\}
\end{array}
\]
\translatornote{等式 \(E_2(\{a\})=\{a\}\) 表示：把集合 \(\{a\}\) 作为输入，
依次检查定义 \(E_2\) 的三条规则。第一条规则没有前提，能够生成 \(a\)；第二条
规则要求输入中有 \(c\)，第三条规则要求输入中同时有 \(a\) 和 \(b\)，因而都
不能使用。最终输出中只有 \(a\)，所以 \(E_2(\{a\})=\{a\}\)。}

\(E_2\)-闭集合为 \(\{a\}\) 和 \(\{a,b,c\}\)；\(E_2\)-一致集合为
\(\varnothing\)、\(\{a\}\) 和 \(\{a,b,c\}\)。所以
\[
\mu E_2=\{a\},\qquad \nu E_2=\{a,b,c\}.
\]

\taplsolutionbacklink{ex:21.1.7}{返回练习21.1.7}
\end{taplauthorsolutionentrybox}

\begin{taplauthorsolutionentrybox}
\phantomsection
\label{sol:author:21.1.9}
\textbf{21.1.9 解答}

为推出自然数上的普通归纳原理，定义
\[
F(X)=\{0\}\cup\{i+1\mid i\in X\}.
\]
把谓词 \(P\) 看成所有满足 \(P\) 的自然数组成的集合。设 \(P(0)\) 成立，
并且 \(P(i)\Rightarrow P(i+1)\)。若 \(X\subseteq P\)，那么 \(F(X)\)
中的 \(0\) 属于 \(P\)；其中每个 \(i+1\) 也因 \(i\in X\subseteq P\) 和
归纳步骤而属于 \(P\)，所以 \(F(X)\subseteq P\)。特别地，取 \(X=P\)，
便有 \(F(P)\subseteq P\)，即 \(P\) 是 \(F\)-闭的。归纳原理给出
\(\mu F\subseteq P\)。另一方面，任何 \(F\)-闭集合都必须包含 \(0\)，并且
只要包含 \(i\) 就必须包含 \(i+1\)，所以 \(\mu F=\mathbb N\)。因此
\(P(n)\) 对所有 \(n\in\mathbb N\) 成立。

对字典序归纳，定义
\[
F(X)=\{(m,n)\mid
  \forall(m',n')<(m,n).\ (m',n')\in X\}.
\]
把谓词 \(P\) 看成所有满足 \(P\) 的自然数对组成的集合。题设是：若某个数对
的所有字典序前驱都属于 \(P\)，那么该数对也属于 \(P\)。因此，若
\(X\subseteq P\)，任取 \((m,n)\in F(X)\)，它的所有字典序前驱都属于
\(X\)，从而也属于 \(P\)；题设于是给出 \((m,n)\in P\)，所以
\(F(X)\subseteq P\)。特别地，\(F(P)\subseteq P\)，即 \(P\) 是
\(F\)-闭的，归纳原理给出 \(\mu F\subseteq P\)。

最后证明 \(\mu F=\mathbb N\times\mathbb N\)。首先，
\(\mathbb N\times\mathbb N\) 是 \(F\)-闭的。再反设它的某个真子集
\(Y\) 也是 \(F\)-闭的。由于字典序是良基的，不属于 \(Y\) 的数对中存在
字典序最小者，记为 \((m,n)\)。它的所有字典序前驱都属于 \(Y\)，所以
\((m,n)\in F(Y)\)；但按选取方式，\((m,n)\notin Y\)，这与
\(F(Y)\subseteq Y\) 矛盾。因此，\(\mathbb N\times\mathbb N\) 的任何
真子集都不可能是 \(F\)-闭的，故最小的 \(F\)-闭集合就是
\(\mathbb N\times\mathbb N\)。

\translatornote{这里定义生成函数 \(F\)，是为了把具体的归纳规则接到
\cref{cor:21.1.8}的抽象归纳原理上。普通归纳中的 \(\{0\}\) 表示起点
“\(0\) 是自然数”，\(\{i+1\mid i\in X\}\) 表示后继规则“已有 \(i\)
便可生成 \(i+1\)”。从空集反复应用 \(F\)，会依次得到 \(\{0\}\)、
\(\{0,1\}\)、\(\{0,1,2\}\) 等集合，最小不动点正是 \(\mathbb N\)。

字典序归纳的生成函数采用同样的思路。条件 \((m,n)\in F(X)\) 表示：所有
字典序小于 \((m,n)\) 的数对都已经属于 \(X\)，因而现在可以处理
\((m,n)\)。例如，若 \(X=\{(0,0)\}\)，则 \((0,1)\) 的唯一字典序前驱是
\((0,0)\)，所以 \((0,1)\in F(X)\)；\((0,0)\) 没有前驱，也属于
\(F(X)\)。但 \((0,2)\) 还有前驱 \((0,1)\notin X\)，因此尚不属于
\(F(X)\)。

证明末尾的“真子集”按集合包含关系比较，“最小数对”则按字典序比较。若
一个真子集 \(Y\) 漏掉的字典序最小数对是 \((m,n)\)，那么它的所有字典序
前驱都已属于 \(Y\)，故 \((m,n)\in F(Y)\)；但 \((m,n)\notin Y\)，所以
\(F(Y)\nsubseteq Y\)，\(Y\) 不可能是 \(F\)-闭的。}

\taplsolutionbacklink{ex:21.1.9}{返回练习21.1.9}
\end{taplauthorsolutionentrybox}

\begin{taplauthorsolutionentrybox}
\phantomsection
\label{sol:author:21.2.2}
\textbf{21.2.2 解答}

先定义一个足够宽的候选全集。这里的树是偏函数
\(T:\{1,2\}^{*}\rightharpoonup\{\to,\times,\tapltype{Top}\}\)，只要求根节点
\(T(\epsilon)\) 有定义，并且只要 \(T(\pi,\sigma)\) 有定义，其父节点
\(T(\pi)\) 就有定义。除此之外暂不限制节点标记；例如，标为
\(\tapltype{Top}\) 的节点也可以有非平凡子节点。把所有这样的树组成的集合
取作全集 \(U\)。与\taplsecref{21.2}{sec:21.2}相同，符号
\(\tapltype{Top}\)、\(\times\) 和 \(\to\) 也表示相应的树构造操作。定义
\[
F(X)=\{\tapltype{Top}\}
\cup\{T_1\times T_2\mid T_1,T_2\in X\}
\cup\{T_1\to T_2\mid T_1,T_2\in X\}
\]

由定义可见 \(\mathcal T\subseteq U\)，所以可以在同一全集中比较
\(\mathcal T\) 与 \(\nu F\)，以及 \(\mathcal T_f\) 与 \(\mu F\)。下面分别
验证这两个等式。

先证 \(\mathcal T=\nu F\)。

\textbf{包含关系 \(\mathcal T\subseteq\nu F\)。}
任取 \(T\in\mathcal T\)。把 \(X=\mathcal T\) 代入 \(F\) 的定义，得到
\[
F(\mathcal T)=\{\tapltype{Top}\}
\cup\{T_1\times T_2\mid T_1,T_2\in\mathcal T\}
\cup\{T_1\to T_2\mid T_1,T_2\in\mathcal T\}
\]
若 \(T=\tapltype{Top}\)，它属于右侧第一部分。若
\(T=T_1\times T_2\) 或 \(T=T_1\to T_2\)，则两棵直接子树仍是合法树，
因而 \(T_1,T_2\in\mathcal T\)，所以 \(T\) 分别属于右侧第二或第三部分。
三种情况下都有 \(T\in F(\mathcal T)\)，故
\(\mathcal T\subseteq F(\mathcal T)\)。这说明 \(\mathcal T\) 是
\(F\)-一致的；再由余归纳原理得到
\(\mathcal T\subseteq\nu F\)。

\textbf{包含关系 \(\nu F\subseteq\mathcal T\)。}
任取 \(T\in\nu F\)。因为 \(\nu F\) 是不动点，所以
\(\nu F=F(\nu F)\)，从而 \(T\in F(\nu F)\)。逐项查看 \(F\) 的定义可知，
树的根节点必由三种构造之一生成：它或者是没有子节点的
\(\tapltype{Top}\)，或者标为 \(\times\) 或 \(\to\)，并有两棵仍属于
\(\nu F\) 的直接子树。对直接子树可以重复同一论证。形式上，沿任意路径
\(\pi\) 对其长度作归纳，便可逐层验证\cref{def:21.2.1}第三、第四项条件。
因此 \(T\) 是合法树类型，即 \(T\in\mathcal T\)。

再证 \(\mathcal T_f=\mu F\)。

\textbf{包含关系 \(\mu F\subseteq\mathcal T_f\)。}
把 \(X=\mathcal T_f\) 代入 \(F\) 的定义：第一部分生成单节点有限树
\(\tapltype{Top}\)，第二、第三部分分别用两棵有限树构造积类型树和箭头
类型树，结果仍是有限树。所以
\(F(\mathcal T_f)\subseteq\mathcal T_f\)，即 \(\mathcal T_f\) 是
\(F\)-闭的。归纳原理于是给出
\(\mu F\subseteq\mathcal T_f\)。

\textbf{包含关系 \(\mathcal T_f\subseteq\mu F\)。}
对有限树的大小作归纳；这里用满足 \(T(\pi)\) 有定义的最长路径 \(\pi\) 的
长度衡量树的大小。单节点树 \(\tapltype{Top}\) 属于 \(F(\mu F)=\mu F\)。
对于更大的有限树，其根标为 \(\times\) 或 \(\to\)，两棵直接子树都更小。
由归纳假设，两棵子树属于 \(\mu F\)；再按 \(F\) 的定义，整棵树也属于
\(F(\mu F)=\mu F\)。因此 \(\mathcal T_f\subseteq\mu F\)。

\taplsolutionbacklink{ex:21.2.2}{返回练习21.2.2}
\end{taplauthorsolutionentrybox}

\begin{taplauthorsolutionentrybox}
\phantomsection
\label{sol:author:21.3.3}
\textbf{21.3.3 解答}

类型对 \((\tapltype{Top},\tapltype{Top}\times\tapltype{Top})\) 不属于
\(\nu\mathcal S\)。从 \(\mathcal S\) 的定义可见，它不属于任何
\(\mathcal S(X)\)，所以不存在包含它的 \(\mathcal S\)-一致集合；特别地，
\(\nu\mathcal S\) 也不包含它。

\taplsolutionbacklink{ex:21.3.3}{返回练习21.3.3}
\end{taplauthorsolutionentrybox}

\begin{taplauthorsolutionentrybox}
\phantomsection
\label{sol:author:21.3.4}
\textbf{21.3.4 解答}

\textbf{无限类型的反例。}
任取无限树类型 \(T\)。下面证明 \((T,T)\in\nu\mathcal S\)，但
\((T,T)\notin\mu\mathcal S\)。令
\[
R=\{(Q,Q)\mid Q\text{ 是 }T\text{ 的子树}\}.
\]
先验证 \(R\subseteq\mathcal S(R)\)。任取 \((Q,Q)\in R\)，按 \(Q\) 的
根节点分类。若 \(Q=\tapltype{Top}\)，则 \((Q,Q)\) 由 \(\mathcal S\) 的
第一部分直接生成。若 \(Q=Q_1\times Q_2\)，则两棵直接子树仍是 \(T\) 的
子树，故 \((Q_1,Q_1),(Q_2,Q_2)\in R\)，积类型分支生成 \((Q,Q)\)。若
\(Q=Q_1\to Q_2\)，箭头类型分支同样由这两个类型对生成 \((Q,Q)\)。因此
\(R\) 是 \(\mathcal S\)-一致的。余归纳给出
\(R\subseteq\nu\mathcal S\)；又因为 \(T\) 是自身的子树，
\((T,T)\in R\)，所以 \((T,T)\in\nu\mathcal S\)。

回到 \((T,T)\)。无限树 \(T\) 至少有一条无限路径。沿这条路径，每遇到
\(\to\) 或 \(\times\) 节点，证明当前子树是自身子类型的推导都必须继续
进入同一棵直接子树，因而形成一条无限推导分支，无法得到有限推导。例如，若
\(T=\tapltype{Top}\times T\)，积类型规则会把 \(T\subtype T\) 化为
\(\tapltype{Top}\subtype\tapltype{Top}\) 和同一个目标
\(T\subtype T\)。所以 \((T,T)\notin\mu\mathcal S\)。

\translatornote{令 \(D\) 为所有具有有限子类型推导的类型对组成的关系。下面
说明 \(D=\mu\mathcal S\)。

先证 \(\mu\mathcal S\subseteq D\)。以 \(D\) 中的判断为前提应用一条规则，
只需把该规则接在若干份有限推导之下，所得推导仍然有限。因此
\(\mathcal S(D)\subseteq D\)，即 \(D\) 是 \(\mathcal S\)-闭的。根据
Knaster--Tarski 定理，\(\mu\mathcal S\) 是所有 \(\mathcal S\)-闭关系的
交，因而包含于每一个闭关系；特别地，\(\mu\mathcal S\subseteq D\)。

再证 \(D\subseteq\mu\mathcal S\)。任取 \(D\) 中的类型对，对其有限推导的
高度作归纳。若最后使用的是无前提规则，其结论直接属于
\(\mathcal S(\mu\mathcal S)=\mu\mathcal S\)。否则，最后一条规则的各个
前提都有更矮的推导，由归纳假设均属于 \(\mu\mathcal S\)；应用最后一条规则，
结论便属于 \(\mathcal S(\mu\mathcal S)=\mu\mathcal S\)。因此
\(D\subseteq\mu\mathcal S\)，最终得到 \(D=\mu\mathcal S\)。所以，只能
依靠无限循环相互支持、无法形成有限推导的判断，不会进入
\(\mu\mathcal S\)。}

\textbf{有限类型的情形。}
不存在属于 \(\nu\mathcal S_f\) 却不属于 \(\mu\mathcal S_f\) 的有限类型对；
事实上，两个不动点相等。按最小不动点的定义，它包含于
\(\mathcal S_f\) 的任意不动点；特别地，
\(\mu\mathcal S_f\subseteq\nu\mathcal S_f\)。故只需证明反向包含关系
\(\nu\mathcal S_f\subseteq\mu\mathcal S_f\)。任取
\((S,T)\in\nu\mathcal S_f\)，对 \(S,T\) 的大小之和归纳。由于
\(\nu\mathcal S_f=\mathcal S_f(\nu\mathcal S_f)\)，可以检查
\(\mathcal S_f\) 中生成该类型对的分支：若 \(T=\tapltype{Top}\)，结论由
无前提规则直接属于 \(\mu\mathcal S_f\)；若 \(S,T\) 都是积类型，两个分量
对应的类型对规模更小，由归纳假设属于 \(\mu\mathcal S_f\)，再应用积类型
分支即可；箭头类型的情形相同，只需注意参数位置的类型对方向相反。于是
\((S,T)\in\mu\mathcal S_f\)，两个不动点相等。

\taplsolutionbacklink{ex:21.3.4}{返回练习21.3.4}
\end{taplauthorsolutionentrybox}

\begin{taplauthorsolutionentrybox}
\phantomsection
\label{sol:author:21.3.8}
\textbf{21.3.8 解答}

定义树类型上的恒等关系
\[
I=\{(T,T)\mid T\in\mathcal T\}.
\]
若能证明 \(I\) 是 \(\mathcal S\)-一致的，余归纳便给出
\(I\subseteq\nu\mathcal S\)，即自反性。任取 \((T,T)\in I\)，对 \(T\) 的
根分类。若 \(T=\tapltype{Top}\)，按定义
\((T,T)\in\mathcal S(I)\)。若 \(T=T_1\times T_2\)，由于
\((T_1,T_1),(T_2,T_2)\in I\)，积类型分支给出
\((T_1\times T_2,T_1\times T_2)\in\mathcal S(I)\)。箭头情形相同。

\taplsolutionbacklink{ex:21.3.8}{返回练习21.3.8}
\end{taplauthorsolutionentrybox}

\begin{taplauthorsolutionentrybox}
\phantomsection
\label{sol:author:21.4.2}
\textbf{21.4.2 解答}

依余归纳原理，只需证明 \(U\times U\) 是
\(F^{\operatorname{TR}}\)-一致的。任取 \((x,y)\in U\times U\)，再任选
\(z\in U\)。有 \((x,z),(z,y)\in U\times U\)，所以按
\(F^{\operatorname{TR}}\) 的定义，
\((x,y)\in F^{\operatorname{TR}}(U\times U)\)。

\taplsolutionbacklink{ex:21.4.2}{返回练习21.4.2}
\end{taplauthorsolutionentrybox}

\Needspace{7\baselineskip}
\begin{taplauthorsolutionentrybox}
\phantomsection
\label{sol:author:21.5.2}
\textbf{21.5.2 解答}

检查可逆性，只需查看 \(\mathcal S_f\) 与 \(\mathcal S\) 的定义，并确认每个
\(\mathcal G_{(S,T)}\) 至多含一个极小成员。两个定义的每一项都明确规定了
可支持元素的形式和支持集内容，所以可直接写出
\(\operatorname{support}_{\mathcal S_f}\) 与
\(\operatorname{support}_{\mathcal S}\)；其形式与
\cref{def:21.8.4}之后的 \(\operatorname{support}_{\mathcal S_\mu}\) 相同，
只去掉两个 \(\mu\) 分支。

\translatornote{把作者省略的分类写开如下。对
\(\mathcal F\in\{\mathcal S_f,\mathcal S\}\)，其支持函数为
\[
\operatorname{support}_{\mathcal F}(S,T)=
\begin{cases}
\varnothing,
  &T=\tapltype{Top},\\
\{(S_1,T_1),(S_2,T_2)\},
  &S=S_1\times S_2,\ T=T_1\times T_2,\\
\{(T_1,S_1),(S_2,T_2)\},
  &S=S_1\to S_2,\ T=T_1\to T_2,\\
\uparrow,
  &\text{其他情形}.
\end{cases}
\]
两者的公式相同；\(\mathcal S_f\) 的参数只取有限树类型，
\(\mathcal S\) 的参数则可以取无限树类型。右侧为
\(\tapltype{Top}\) 时无需任何前提，所以最小支持集为空集。积类型必须同时
支持两个对应分量；箭头类型必须支持逆向的参数类型对
\((T_1,S_1)\) 和正向的结果类型对 \((S_2,T_2)\)。其余外层形状无法由任何
分支生成，故支持函数无定义。这几种形状互不重叠，而且每种形状所需的前提集合
都唯一，因此 \(\mathcal S_f\) 与 \(\mathcal S\) 都可逆。}

\taplsolutionbacklink{ex:21.5.2}{返回练习21.5.2}
\end{taplauthorsolutionentrybox}

\begin{taplauthorsolutionentrybox}
\phantomsection
\label{sol:author:21.5.4}
\textbf{21.5.4 解答}

相应推导规则是
\[
\frac{i\quad a}{h}
\qquad
\frac{b\quad c}{a}
\qquad
\frac{b}{d}
\qquad
\frac{d}{e}
\qquad
\frac{e}{b}
\qquad
\frac{f\quad g}{c}
\qquad
\frac{g}{f}
\qquad
\frac{}{g}.
\]
由这些规则可直接计算题目给出的两个集合；图中的
\(\mu E=\{c,f,g\}\) 与
\(\nu E=\{a,b,c,d,e,f,g\}\) 也分别满足最小、最大不动点的定义。

\taplsolutionbacklink{ex:21.5.4}{返回练习21.5.4}
\end{taplauthorsolutionentrybox}

\begin{taplauthorsolutionentrybox}
\phantomsection
\label{sol:author:21.5.6}
\textbf{21.5.6 解答}

逆命题不成立。\(\nu F\setminus\mu F\) 中的元素也可能通向一条无限链，而非
环。例如定义
\[
F:\powerset(\mathbb N)\to\powerset(\mathbb N),\qquad
F(X)=\{0\}\cup\{n\mid n+1\in X\}.
\]
则 \(\mu F=\{0\}\)，\(\nu F=\mathbb N\)。对每个 \(n>0\)，有
\(\operatorname{support}(n)=\{n+1\}\)，所以从 \(n\) 出发得到没有环的无限链。

\taplsolutionbacklink{ex:21.5.6}{返回练习21.5.6}
\end{taplauthorsolutionentrybox}

\begin{taplauthorsolutionentrybox}
\phantomsection
\label{sol:author:21.5.13}
\textbf{21.5.13 解答}

先证部分正确性。两项都对一次算法运行的递归结构作归纳。
\begin{enumerate}
  \item 若算法因 \(X=\varnothing\) 返回真，则 \(X\subseteq\mu F\) 显然。
        若因为 \(\operatorname{lfp}(\operatorname{support}(X))\) 返回真，归纳
        假设给出 \(\operatorname{support}(X)\subseteq\mu F\)，再由
        \cref{lem:21.5.8}得到 \(X\subseteq\mu F\)。
  \item 若因 \(\operatorname{support}(X)\uparrow\) 返回假，
        \cref{lem:21.5.8}给出 \(X\not\subseteq\mu F\)。否则递归调用返回假，
        归纳假设给出 \(\operatorname{support}(X)\not\subseteq\mu F\)，再由
        同一引理得到结论。
\end{enumerate}

下面刻画保证算法对有限输入终止的一类生成函数。对有限状态生成函数 \(F\)，令
偏函数 \(\operatorname{height}_F:U\rightharpoonup\mathbb N\) 为满足以下条件的
最小偏函数：\footnote{也可以把这个定义改写为一个最小不动点：在表示偏函数的
关系上定义适当的单调函数，再取其最小不动点。}
\[
\operatorname{height}(x)=
\begin{cases}
0,&\operatorname{support}(x)=\varnothing,\\
0,&\operatorname{support}(x)\uparrow,\\
1+\max\{\operatorname{height}(y)\mid y\in\operatorname{support}(x)\},
  &\operatorname{support}(x)\ne\varnothing.
\end{cases}
\]
若 \(x\) 位于可达环中，或依赖环中元素，则高度无定义。若
\(\operatorname{height}_F\) 是全函数，称 \(F\) 具有有限高度。若
\(y\in\operatorname{support}(x)\) 且二者高度都有定义，则
\(\operatorname{height}(y)<\operatorname{height}(x)\)。

若 \(F\) 是有限状态的，并且具有有限高度，则 \(\operatorname{lfp}(X)\) 对
每个有限输入都终止。递归调用中的集合始终有限，而
\[
h(Y)=\max\{\operatorname{height}(y)\mid y\in Y\}
\]
每次严格减小且非负，可作为终止度量。

\taplsolutionbacklink{ex:21.5.13}{返回练习21.5.13}
\end{taplauthorsolutionentrybox}

\begin{taplauthorsolutionentrybox}
\phantomsection
\label{sol:author:21.8.5}
\textbf{21.8.5 解答}

\(\mathcal S_d\) 与 \(\mathcal S_\mu\) 的定义相同，只是最后一项不要求
\(T\ne\mu X.T_1\) 和 \(T\ne\tapltype{Top}\)。它不可逆，因为
\(\mathcal G_{(\mu X.\tapltype{Top},\mu Y.\tapltype{Top})}\) 含有两个彼此不
包含的极小生成集
\[
\{(\tapltype{Top},\mu Y.\tapltype{Top})\},\qquad
\{(\mu X.\tapltype{Top},\tapltype{Top})\}.
\]
可将这里的两个极小生成集，与同一类型对在 \(\mathcal S_\mu\) 下的生成集
作比较。

\(\mathcal S_\mu\) 的最后一项是 \(\mathcal S_d\) 最后一项的限制，其他各项
相同，因此
\(\nu\mathcal S_\mu\subseteq\nu\mathcal S_d\)。反向包含关系由余归纳和下面
的引理得到；该引理说明 \(\nu\mathcal S_d\) 是
\(\mathcal S_\mu\)-一致的。

\setcounter{taplappendixlemma}{16}
\begin{taplappendixlemma}
\label{lem:a.17}
对任意 \(\mu\)-类型 \(S,T\)，若
\((S,T)\in\nu\mathcal S_d\)，则
\((S,T)\in\mathcal S_\mu(\nu\mathcal S_d)\)。
\end{taplappendixlemma}

\begin{proof}[证明概要]
令 \(k=\mu\text{-}\operatorname{height}(S)\)，对 \(k\) 作归纳。
归纳把
\((S,T)\in\nu\mathcal S_d\) 的任意推导变换为同一判断在
\(\nu\mathcal S_\mu\) 中的推导。左 \(\mu\)-折叠规则的限制保证，变换后的每段
折叠序列都先连续应用若干次左折叠，再连续应用若干次右折叠。
\end{proof}

\translatornote{这项重排是要把 \(\mathcal S_d\) 中允许任意穿插的折叠步骤，
改造成符合 \(\mathcal S_\mu\) 限制的推导。按推导树从前提向结论读，一旦右
折叠在右侧加入最外层 \(\mu\)，\(\mathcal S_\mu\) 左折叠分支要求“右侧不以
\(\mu\) 开头”的旁条件便不再满足；因此每段折叠只能先做左折叠，再做右折叠。
重排的难点只来自这项左折叠限制；右折叠没有相应的旁条件。沿推导树从结论
追溯前提时，每处理一次左折叠，左侧类型开头连续出现的 \(\mu\) 绑定便少一个，
所以只需对这个数量 \(k\) 归纳。}

\taplsolutionbacklink{ex:21.8.5}{返回练习21.8.5}
\end{taplauthorsolutionentrybox}

\begin{taplauthorsolutionentrybox}
\phantomsection
\label{sol:author:21.9.2}
\textbf{21.9.2 解答}

\[
\frac{}{T\sqsubseteq T}
\qquad
\frac{S\sqsubseteq T_1}{S\sqsubseteq T_1\times T_2}
\qquad
\frac{S\sqsubseteq T_2}{S\sqsubseteq T_1\times T_2}
\]
\[
\frac{S\sqsubseteq T_1}{S\sqsubseteq T_1\to T_2}
\qquad
\frac{S\sqsubseteq T_2}{S\sqsubseteq T_1\to T_2}
\qquad
\frac{S\sqsubseteq[X\mapsto\mu X.T]T}{S\sqsubseteq\mu X.T}.
\]

有一点值得注意：\(\operatorname{TD}\) 与本章此前考察的生成函数不同，
它不可逆。例如判断
\(B\sqsubseteq(A\times B)\to(B\times C)\) 有两个生成集
\(\{B\sqsubseteq A\times B\}\) 与
\(\{B\sqsubseteq B\times C\}\)，二者互不包含。

\taplsolutionbacklink{ex:21.9.2}{返回练习21.9.2}
\end{taplauthorsolutionentrybox}

\begin{taplauthorsolutionentrybox}
\phantomsection
\label{sol:author:21.9.7}
\textbf{21.9.7 解答}

除以 \(\mu\) 开头的类型外，\(\operatorname{BU}\) 的规则与
\hyperref[sol:author:21.9.2]{21.9.2 解答中的 \(\operatorname{TD}\) 规则}相同；
\(\mu\) 分支改为
\[
\frac{S\preccurlyeq T}
     {[X\mapsto\mu X.T]S\preccurlyeq\mu X.T}.
\]

\taplsolutionbacklink{ex:21.9.7}{返回练习21.9.7}
\end{taplauthorsolutionentrybox}

\begin{taplauthorsolutionentrybox}
\phantomsection
\label{sol:author:21.11.1}
\textbf{21.11.1 解答}

例子很多。对几乎任意 \(T\)，\(\mu X.T\) 与
\([X\mapsto\mu X.T]T\) 就是一个简单例子。更有意思的例子是
\[
\mu X.\tapltype{Nat}\times(\tapltype{Nat}\times X)
\qquad\text{和}\qquad
\mu X.\tapltype{Nat}\times X.
\]

\taplsolutionbacklink{ex:21.11.1}{返回练习21.11.1}
\end{taplauthorsolutionentrybox}
